% ═══════════════════════════════════════════════════════════════════════════════
\chapter{Results \& Discussion}
% ═══════════════════════════════════════════════════════════════════════════════

% overview
% remember to include reference to the other simualtion papers throughout to validate the results.

\section{Baseline Protocol Verification}
\label{sec:baseline}

To establish a control baseline, the protocols are initially evaluated over an ideal \gls{qchan}, absent of environmental noise or adversarial interference. This first experiment compares the operational metrics of \gls{bb84}, \gls{b92}, and \gls{e91} to verify their software implementations. As depicted in \Cref{fig:baseline}, the mutual information for all three protocols converges to the theoretical maximum. 

% mutual info bar chart
\begin{figure}[htbp]
    \centering
    \includegraphics[width=0.5\linewidth]{images.results v2/exp1/exp1_baseline_mutual_info.png}
    \caption{Baseline Mutual Information under Ideal Conditions.}
    \label{fig:baseline}
\end{figure}

\Cref{tab:validation} presents the theoretically expected values alongside the simulated means, where every simulated value is expected to lie within $\pm 2\%$ of its theoretical target. This outcome confirms that the simulation architecture processes quantum states without degradation. As such, it validates that the system provides a sound control environment prior to introducing interference models, aligning with the approach of recent \gls{qkd} simulation platforms \cite{devi2025}.

\begin{table}[htbp]
    \centering
    \caption{Ideal Channel Validation Results ($n=1000$, $N_{trials}=50$)}
    \label{tab:validation}
    \footnotesize
    \begin{tabular}{@{}l c c@{}}
        \toprule
        \textbf{Metric} & \textbf{Target} & \textbf{Simulated} \\
        \midrule
        \gls{bb84} sifting efficiency (\%) & $50.0$ & $50.0$ \\
        \gls{b92} sifting efficiency (\%)  & $25.0$ & $25.0$ \\
        \gls{e91} sifting efficiency (\%)  & $22.2$ & $22.2$ \\
        \gls{bb84} ideal-channel \gls{qber} (\%) & $0.0$  & $0.0$  \\
        \gls{b92} ideal-channel \gls{qber} (\%)  & $0.0$  & $0.0$  \\
        \gls{e91} ideal-channel \gls{qber} (\%)  & $0.0$  & $0.0$  \\
        \gls{e91} \gls{chsh} parameter \(|S|\)   & \(2\sqrt{2}\!\approx\!2.828\) & $2.828$ \\
        \bottomrule
    \end{tabular}
\end{table}


\section{Resilience to Channel Noise}
\label{sec:resilience}

Practical \gls{qkd} deployment requires resilience against environmental decoherence. The depolarising channel model introduces random state errors, which inherently degrade the transmitted quantum states. As modelled in \Cref{fig:qbernoise}, increasing Noise Strength directly elevates the \gls{qber} across all three protocols.

% qber vs noise - increase
\begin{figure}[htbp]
    \centering
    \includegraphics[width=0.5\linewidth]{images.results v2/exp2/exp2_qber_vs_noise.png}
    \caption{QBER vs. Noise Strength, under the Depolarising Channel Model.}
    \label{fig:qbernoise}
\end{figure}


The cascading impact of this elevated \gls{qber} on the system's information-theoretic security must be acknowledged. \Cref{fig:mutualnoise} illustrates how the rising \gls{qber} degrades the mutual information shared between Alice and Bob, as the channel's capacity to carry uncorrupted quantum data diminishes.

% mutual info vs noise - degredation
\begin{figure}[htbp]
    \centering
    \includegraphics[width=0.5\linewidth]{images.results v2/exp2/exp2_mutual_info_vs_noise.png}
    \caption{Mutual Information vs. Noise Strength, under the Depolarising Channel Model.}
    \label{fig:mutualnoise}
\end{figure}

Ultimately, this reduction in mutual information fundamentally restricts the final cryptographic yield. This behaviour is depicted in \Cref{fig:keyratenoise}, which plots the \gls{skr} against Noise Strength for \gls{bb84}, \gls{b92}, and E91. The comparative analysis clearly demonstrates that at a channel noise strength of $5\%$, \gls{bb84} sustains a \gls{skr} of $\sim24\ \text{bits/100 qubits}$, whereas \gls{e91}'s \gls{skr} reduces to $\sim7\ \text{bits/100 qubits}$, and \gls{b92}'s \gls{skr} collapses to $\sim3\ \text{bits/100 qubits}$.

% secure key rates vs noise - degredation
\begin{figure}[htbp]
    \centering
    \includegraphics[width=0.5\linewidth]{images.results v2/exp2/exp2_secure_key_rate_vs_noise.png}
    \caption{Secure Key Rate vs. Noise Strength, under the Depolarising Channel Model.}
    \label{fig:keyratenoise}
\end{figure}

The steeper degradation curve of \gls{b92} stems from its reliance on only two non-orthogonal states. Consequently, depolarising noise rapidly introduces measurement errors that Bob cannot resolve, forcing a severe discard rate during sifting \cite{bennett1992}. Interestingly, despite consistently underperforming \gls{b92} in raw mutual information, \gls{e91} demonstrates superior noise resilience by sustaining a greater \gls{skr} at noise probabilities above $P \geq 2.5\%$. While \gls{b92} rapidly approaches its Devetak-Winter threshold, \gls{e91} leverages its entanglement-based correlations to maintain a functional but reduced cryptographic key rate.


\section{Eavesdropper Vulnerability}
\label{sec:everesults}

The fundamental security guarantee of quantum cryptography relies on the measurement-disturbance principle. To isolate and quantify this vulnerability, this experiment evaluates the \gls{qber} performance of \gls{bb84} and \gls{b92} under an active intercept-resend attack, operating over an ideal QC. The resultant \gls{qber} against Eve's interception rate is depicted in \Cref{fig:qberEve}.

% \gls{qber} increase 2D plot
\begin{figure}[htbp]
    \centering
    \includegraphics[width=0.5\linewidth]{images.results v2/exp3/exp3_eve_vulnerability_qber.png}
    \caption{QBER vs. Intercept Rate, under an Ideal Channel Model.}
    \label{fig:qberEve}
\end{figure}

Crucially, the graph incorporates horizontal threshold lines representing the Devetak-Winter bounds, identifying the point where \gls{skr} drops to zero. For \gls{bb84}, the analysis confirms the widely established $11\%$ \gls{qber} security limit \cite{ShorPreskill2000}. The plotted data demonstrates that this critical threshold is breached at an Intercept Rate of precisely 0.43. Beyond this threshold, the protocol must abort.

In contrast, the \gls{b92} protocol exhibits much higher sensitivity to interception. The simulation records \gls{b92} reaching its theoretical \gls{qber} threshold of $6.5\%$ at a mere 0.13 Intercept Rate, as indicated by the perpendicular axis lines. These results validate that while \gls{b92} offers a simplified architecture, its tolerance against eavesdropping is severely compromised.

\section{E91 Entanglement Mechanics}
\label{sec:e91results}

To verify the entanglement mechanics established in \Cref{subsec:CHSH}, this experiment models the \gls{e91} protocol under a depolarising channel. \Cref{fig:chshqber} presents a dual-axis evaluation, adapted from the simulation framework by Rasay \textit{et al.} \cite{rasay2026}, plotting the simulated S-parameter and \gls{qber} against increasing Noise Strength.

% \gls{chsh} \gls{qber} 2 axis plot
\begin{figure}[htbp]
    \centering
    \includegraphics[width=0.5\linewidth]{images.results v2/exp4/exp4_chsh_qber.png}
    \caption{Dual-Axis \gls{chsh} \& \gls{qber} vs. Noise Strength, under the depolarising channel model.}
    \label{fig:chshqber}
\end{figure}

This data maps the channel's transition from a quantum to a classical state, with the graph presenting the Noise Strength threshold of $P\approx0.165$, where interference disrupts the photon pairs, degrading the S-parameter below the classical limit of $2$. Interestingly, at the breach of this limit, $\text{QBER}\approx17\%$ enabling the inference of an approximate \gls{qber} threshold for this particular \gls{e91} experiment configuration. Below this classical limit, non-local quantum correlations cannot be guaranteed. To quantify the cryptographic penalty of this degradation, \Cref{fig:e91keyrate} maps the resulting \gls{skr} against the raw mutual information.

% secure key rate vs noise, overlaid with mutual info to highlight the separation between mutual info and usable key rate
\begin{figure}[htbp]
    \centering
    \includegraphics[width=0.5\linewidth]{images.results v2/exp4/exp4_secure_key_rate.png}
    \caption{E91 Secure Key Rate vs. Noise, under the depolarising channel model, with Mutual Information overlaid.}
    \label{fig:e91keyrate}
\end{figure}

This graph illustrates the divergence between channel capacity and usable key generation. As \gls{qber} rises toward the threshold, the overhead required for key generation accelerates non-linearly. Thus, the \gls{skr} reaches zero significantly earlier than the mutual information curve, confirming that a noisy channel retains residual classical information capacity long after its ability to produce a secure key has collapsed.


\section{Multivariable System Dynamics}
\label{sec:multi}

Considering a real-world quantum channel, interference would rarely originate from an isolated source, typically suffering simultaneous degradation from environmental noise and adversarial interception. To evaluate these multivariable dynamics, these experiments present 3D surface plots for \gls{bb84} and B92. These final models map Noise Strength and Eve's Intercept Rate against the resultant \gls{qber}, using theoretical threshold planes to visually depict the volumetric boundaries where \gls{skr} drops to zero, rendering key distillation impossible.

\subsection{BB84 Performance}
\label{subsec:bb84perf}

\subsubsection{QBER Surface}
\label{subsubsec:bb84QBERsurface}

As previously described in \ref{sec:protocolimplementation}, the security degradation of \gls{pandm} is evaluated by the \gls{qber} of a protocol. Hence, to evaluate the performance of \gls{bb84}, the \gls{qber} surface modelled in \Cref{fig:bb84QBERsurface} depicts operation, bounded by the well-established $11\%$ \gls{qber} threshold plane.

% \gls{qber} surface plot
\begin{figure}[htbp]
    \centering
    \includegraphics[width=0.5\linewidth]{images.results v2/exp5/exp5_bb84_qber_v4.png}
    \caption{BB84 Multivariable \gls{qber} Surface.}
    \label{fig:bb84QBERsurface}
\end{figure}

Analysing the \gls{qber} surface gradient reveals how environmental noise and adversarial interception interact, compressing the protocol's security margins. \Cref{tab:bb84QBERpoints} details these interference permutations through surface sampling. This data demonstrates that even a minimal baseline of channel noise accelerates the \gls{qber} trajectory, significantly reducing the maximum tolerable intercept rate before the $11\%$ plane is breached. However, \gls{bb84} still maintains a significant "secure operational volume" beneath the \gls{qber} threshold, validating key distribution in moderately noisy environments.

% \gls{qber} surface
% match data for the 6 points, 2 secure, 4 insecure
\begin{table}[htbp]
    \centering
    \caption{BB84 \gls{qber} Surface Samples ($n_{qubits}=500$, $N_{trials}=20$).}
    \label{tab:bb84QBERpoints}
    \footnotesize
    \begin{tabular}{@{}l c c c@{}}
        \toprule
        \textbf{Eve Rate (\%)} & \textbf{Noise Strength (\%)} & \textbf{QBER} & \textbf{Security}\\
        \midrule
        $7.50\%$ & $8.00\%$ & \textbf{$5.43\%$} & \textbf{Secure}  \\
        $22.5\%$ & $4.00\%$ & \textbf{$7.26\%$} & \textbf{Secure}  \\
        $6.00\%$ & $20.0\%$ & \textbf{$11.3\%$} & Insecure  \\
        $30.0\%$ & $13.0\%$ & \textbf{$12.9\%$} & Insecure  \\
        $39.0\%$ & $21.0\%$ & \textbf{$18.5\%$} & Insecure  \\
        $46.5\%$ & $24.0\%$ & \textbf{$21.3\%$} & Insecure  \\
        \bottomrule
    \end{tabular}
\end{table}

\subsubsection{Mutual Information Surface}
\label{subsubsec:bb84MIsurface}
Prior to analysing \gls{bb84}'s \gls{skr} performance, it is critical to comment upon the relationship between the Mutual Information shared by Alice and Bob, and the Mutual Information shared between Alice and Eve. The surface model in \Cref{fig:bb84MIsurface} depicts this behaviour using two surfaces, each to represent their respective Mutual Informations.

% MI surface plot
\begin{figure}[htbp]
    \centering
    \includegraphics[width=0.5\linewidth]{images.results v2/exp5/exp5_bb84_mutual_info_v4.png}
    \caption{BB84 Multivariable Mutual Information Surfaces.}
    \label{fig:bb84MIsurface}
\end{figure}


Reflecting the relationship between Mutual Information during channel interference described in \ref{subsec:mutualinfo}, the data in \Cref{tab:bb84MIsurface} presents sample Mutual Information rates across the surface.

% MI surface
\begin{table}[htbp]
    \centering
    \caption{BB84 Mutual Information Surface Samples ($n_{qubits}=500$, $N_{trials}=20$).}
    \label{tab:bb84MIsurface}
    \footnotesize
    \begin{tabular}{@{}l c c c@{}}
        \toprule
        \textbf{Eve Rate (\%)} & \textbf{Noise Strength (\%)} & \textbf{Mutual Information (bits per 100 qubits)} & \textbf{Security}\\
        \midrule
        $4.5\%$ & $4.00\%$ & \textbf{$39.5$} & \textbf{Secure}  \\
        $24.0\%$ & $3.00\%$ & \textbf{$31.4$} & \textbf{Secure}  \\
        $3.00\%$ & $16.0\%$ & \textbf{$28.6$} & \textbf{Secure}  \\
        $10.5\%$ & $22.0\%$ & \textbf{$22.6$} & Insecure  \\
        $46.5\%$ & $20.0\%$ & \textbf{$15.2$} & Insecure  \\
        $49.5\%$ & $26.0\%$ & \textbf{$11.9$} & Insecure  \\
        \bottomrule
    \end{tabular}
\end{table}


\subsubsection{Secure Key Rate Surface}
\label{subsubsec:bb84SKRsurface}

Ultimately, the key metric that defines the utility of a \gls{qkd} protocol is the \gls{skr}, which dictates the ability to perform key distribution. As such, the \gls{skr} performance of \gls{bb84} under multivariable interference is modelled by the surface in \Cref{fig:bb84SKRsurface}.

% \gls{skr} surface plot
\begin{figure}[htbp]
    \centering
    \includegraphics[width=0.5\linewidth]{images.results v2/exp5/exp5_bb84_skr_v4.png}
    \caption{BB84 Multivariable Secure Key Rate Surface.}
    \label{fig:bb84SKRsurface}
\end{figure}

Identifying a sample of \gls{skr} performance under the aforementioned channel conditions, the data in \Cref{tab:bb84KRsurface} details interference permutations across the surface.

% \gls{skr} surface
\begin{table}[htbp]
    \centering
    \caption{BB84 Secure Key Rate Surface Samples ($n_{qubits}=500$, $N_{trials}=20$).}
    \label{tab:bb84KRsurface}
    \footnotesize
    \begin{tabular}{@{}l c c@{}}
        \toprule
        \textbf{Eve Rate (\%)} & \textbf{Noise Strength (\%)} & \textbf{SKR (bits per 100 qubits)}\\
        \midrule
        $6.00\%$ & $4.00\%$ & \textbf{$27.8$} \\
        $4.50\%$ & $9.00\%$ & \textbf{$18.0$} \\
        $12.0\%$ & $13.0\%$ & \textbf{$6.12$} \\
        $28.5\%$ & $6.00\%$ & \textbf{$4.48$} \\
        $6.00\%$ & $19.0\%$ & \textbf{$2.22$} \\
        $36.0\%$ & $17.0\%$ & \textbf{$0.00$} \\    
        \bottomrule
    \end{tabular}
\end{table}

%%%%%%%%%%%%%%%%%%                           b92                                %%%%%%%%%%%%%%%%%%%
\subsection{B92 Performance}
\label{subsec:b92perf}

\subsubsection{QBER Surface}
\label{subsubsec:b92QBERsurface}

Looking to \Cref{fig:b92QBERsurface}, the plot evaluates \gls{b92} against its stricter $6.5\%$ security threshold. The rapid ascent of the \gls{qber} surface along both axes further demonstrates \gls{b92}'s fragility under channel interference. Due to its two-state nature, the presence of both channel depolarisation and eavesdropping triggers a disproportionate increase in errors in comparison to BB84. As a result, the "secure operational volume" beneath the threshold plane is highly restricted.

% \gls{qber} surface plot
\begin{figure}[htbp]
    \centering
    \includegraphics[width=0.5\linewidth]{images.results v2/exp6/exp6_b92_qber_v4.png}
    \caption{B92 Multivariable \gls{qber} Surface.}
    \label{fig:b92QBERsurface}
\end{figure}

Analysing the \gls{qber} surface gradient reveals how combined interference drastically compresses \gls{b92}'s tolerance to noise. \Cref{tab:b92QBERpoints} details a sample of interference permutations, demonstrating that, despite its fragility, \gls{b92} is still able to securely distribute cryptographic material, albeit under a reduced noise regime compared to BB84.

% \gls{qber} surface
% match data for the 6 points, 2 secure, 4 insecure
\begin{table}[htbp]
    \centering
    \caption{B92 \gls{qber} Surface Samples ($n_{qubits}=500$, $N_{trials}=20$).}
    \label{tab:b92QBERpoints}
    \footnotesize
    \begin{tabular}{@{}l c c c@{}}
        \toprule
        \textbf{Eve Rate (\%)} & \textbf{Noise Strength (\%)} & \textbf{QBER} & \textbf{Security}\\
        \midrule
        $2.25\%$ & $2.00\%$ & \textbf{$3.32\%$} & \textbf{Secure}  \\
        $9.00\%$ & $1.50\%$ & \textbf{$5.37\%$} & \textbf{Secure}  \\
        $1.50\%$ & $5.50\%$ & \textbf{$5.60\%$} & Insecure  \\
        $9.75\%$ & $3.50\%$ & \textbf{$7.84\%$} & Insecure  \\
        $6.75\%$ & $6.50\%$ & \textbf{$9.14\%$} & Insecure  \\
        $3.75\%$ & $11.0\%$ & \textbf{$10.8\%$} & Insecure  \\
        \bottomrule
    \end{tabular}
\end{table}

\subsubsection{Mutual Information Surface}
\label{subsubsec:b92MIsurface}
Once again, establishing the Mutual Information zones for \gls{b92} is critical to inform the analysis of subsequent \gls{skr} performance. \Cref{fig:b92MIsurface} models this relationship by plotting the respective Mutual Information surfaces on the same graph, producing a visual tool to understand how Mutual Information, in turn, informs \gls{skr} calculations. 

% MI surface plot
\begin{figure}[htbp]
    \centering
    \includegraphics[width=0.5\linewidth]{images.results v2/exp6/exp6_b92_mutual_info_v4.png}
    \caption{B92 Multivariable Mutual Information Surfaces.}
    \label{fig:b92MIsurface}
\end{figure}

Once more, reflecting the relationship between Mutual Information during channel interference described in \ref{subsec:mutualinfo}, the data in \Cref{tab:b92MIsurface} synthesises sample Mutual Information values across the surface. Interestingly, beyond the point at which the channel is considered insecure, Mutual Information between Alice and Bob still exists, albeit exceeded by the Mutual Information of Alice and Eve, resulting in a compromised system.

% MI surface
\begin{table}[htbp]
    \centering
    \caption{B92 Mutual Information Surface Samples ($n_{qubits}=500$, $N_{trials}=20$).}
    \label{tab:b92MIsurface}
    \footnotesize
    \begin{tabular}{@{}l c c c@{}}
        \toprule
        \textbf{Eve Rate (\%)} & \textbf{Noise Strength (\%)} & \textbf{Mutual Information (bits per 100 qubits)} & \textbf{Security}\\
        \midrule
        $1.50\%$ & $1.00\%$ & \textbf{$22.1$} & \textbf{Secure}  \\
        $2.25\%$ & $3.00\%$ & \textbf{$20.3$} & \textbf{Secure}  \\
        $16.5\%$ & $2.00\%$ & \textbf{$13.9$} & Insecure  \\
        $13.5\%$ & $7.00\%$ & \textbf{$12.2$} & Insecure  \\
        $10.5\%$ & $12.5\%$ & \textbf{$11.9$} & Insecure  \\ 
        $24.0\%$ & $9.50\%$ & \textbf{$8.60$} & Insecure  \\ 
        \bottomrule
    \end{tabular}
\end{table}


\subsubsection{Secure Key Rate Surface}
\label{subsubsec:b92SKRsurface}
Described in \ref{subsubsec:bb84SKRsurface}, \gls{skr} is critical in determining the rate at which cryptographic material can be securely distributed by a protocol. Acknowledging this, \Cref{fig:b92SKRsurface} models the \gls{skr} performance through a resultant \gls{skr} surface under multivariable interference.

% skr surface plot
\begin{figure}[htbp]
    \centering
    \includegraphics[width=0.5\linewidth]{images.results v2/exp6/exp6_b92_skr_v4.png}
    \caption{B92 Multivariable Secure Key Rate Surface.}
    \label{fig:b92SKRsurface}
\end{figure}

Examining the \gls{skr} performance during channel interference, the data in \Cref{tab:b92SKRsurface} exhibits sample Secure Key Rates across the surface. Crucially, these multivariable results empirically prove that \gls{b92} requires near-perfect channel conditions to remain cryptographically secure, reducing its deployment viability in real-world applications.

% \gls{skr} surface
\begin{table}[htbp]
    \centering
    \caption{B92 Secure Key Rate Surface Samples ($n_{qubits}=500$, $N_{trials}=20$).}
    \label{tab:b92SKRsurface}
    \footnotesize
    \begin{tabular}{@{}l c c@{}}
        \toprule
        \textbf{Eve Rate (\%)} & \textbf{Noise Strength (\%)} & \textbf{SKR (bits per 100 qubits)}\\
        \midrule
        $1.50\%$ & $0.50\%$ & \textbf{$17.9$} \\
        $1.50\%$ & $2.00\%$ & \textbf{$12.3$} \\
        $5.25\%$ & $1.50\%$ & \textbf{$7.25$} \\
        $5.25\%$ & $3.00\%$ & \textbf{$3.25$} \\
        $9.75\%$ & $1.00\%$ & \textbf{$1.93$} \\
        $21.0\%$ & $7.00\%$ & \textbf{$0.00$} \\
        \bottomrule
    \end{tabular}
\end{table}

It must be said that whilst these surface graphs depict the convolved effects of multiple interference sources, in practical deployment, distinguishing between \gls{qchan} noise and an eavesdropper attack is notoriously difficult, even with an established operational noise regime known to both parties prior to transmission. 
Specifically for \gls{skr} calculations, as described in \Cref{subsec:DWsec}, a worst case scenario scheme must be employed, attributing all noise to the presence of an eavesdropper. Despite this, they provide a useful visual aid for understanding the multivariable system dynamics of \gls{bb84} and B92.
